\section{Evaluation protocol}
\label{sec:data}

All results are on AMLworld HI-Small (5{,}078{,}345 transactions) under the IBM day-based
60/20/20 split: train days 0--5, validation days 6--7, test days 8--17.
% source: docs/report.md §3.3
Unless stated otherwise, numbers are over 5 detector seeds, on the \emph{primary window}
(test days 8--9, 174 in-scope truth rings), under the \texttt{test\_visible} scope at
Jaccard $\ge 0.5$ matching.
% source: docs/FINDINGS.md (preamble)

\subsection{Ring ground truth and scope}
\label{sec:data-scope}

A ground-truth ring is one laundering-attempt block in the generator's pattern file; HI-Small
has 370 rings over 3{,}209 pattern transactions in eight topologies (FAN-OUT, FAN-IN,
GATHER-SCATTER, SCATTER-GATHER, CYCLE, RANDOM, BIPARTITE, STACK).
% source: docs/RING_GROUND_TRUTH_DEFINITION.md §1–2
A ring's members are the accounts appearing in its transactions, and a prediction is scored
against that account set.

A ring's transactions can straddle split boundaries, so the scope rule --- which rings a
window is held responsible for, and with which members --- must be fixed and stated. We use
\texttt{test\_visible}: a ring is in scope for a window if it has at least one transaction
there, and its member set is the accounts of those in-window transactions. Two stricter
alternatives (every transaction in the window; more than half) are reported in the released
tables and never used for a headline.
% source: docs/paper/evaluation_protocol.md §2
Even under this scope, membership is only partly observable in one split: 7.9\,\% of ring
members have no test-period transaction at all, and 33.2\,\% have none labelled illicit.
% source: docs/FINDINGS.md (C5: 155/1,962 and 652/1,962)
Section~\ref{sec:findings-scope} measures what the scope choice is worth.

\subsection{Matching rule}

A predicted ring matches a truth ring when their account sets reach Jaccard $\ge 0.5$, paired
greedily one-to-one by descending Jaccard. Two truth rings can both reach Jaccard $\ge 0.5$ with one
prediction (at exactly 0.5), so ties are possible; they are broken by the smaller truth-ring
identifier, then by the prediction's sorted member list, which makes the result independent of file
order. Against the file-order tie-break used for the reported tables, no match count changes in any
headline cell, and the largest change in a five-seed mean ring F1 in any cell is 0.0032.
% source: results/notes/matcher_tiebreak.md; docs/PROTOCOL.md (Matching)
% source: docs/paper/evaluation_protocol.md §3
The threshold is fixed, never tuned, and results at 0.3 and 0.7 are reported beside every
headline, because a ring-level gain quoted at one threshold can be a fragmentation artefact
(Section~\ref{sec:results-partition} gives an instance from our own pipeline). Predictions
smaller than three accounts are discarded (\texttt{min\_size} $=3$). This constrains the
unrestricted prediction; matching is done on predictions restricted to the evaluation window,
which can be smaller, so \texttt{min\_size} does not imply a ceiling on ring recall
(Section~\ref{sec:results-ceiling}).
% source: docs/paper/evaluation_protocol.md §3

\subsection{Evaluation windows}
\label{sec:data-windows}

Every headline uses the \emph{primary} window (test days 8--9); the tail (days 10--17) and the
whole test period are reported separately (Section~\ref{sec:findings-tail} gives the reason).
% source: docs/paper/evaluation_protocol.md §4, §9; window details in docs/PROTOCOL.md

\subsection{Nothing is chosen on the test period}

The seed fraction is fixed in advance at $k^{*}=0.1\,\%$. Every other parameter --- \texttt{min\_size}, the expansion $\Delta t$
and probability floor, the linkage cut, the calibration map, the detector checkpoint --- is
selected on the validation split and frozen; the runners accept no other choice.
% source: docs/paper/evaluation_protocol.md §5
Two consequences cost us numbers and are reported anyway: the validation-chosen operating
point is not test-optimal (released repository), and when a pattern type is held out
its labels are removed from train \emph{and} validation together, because validation chooses
every parameter. Calibration is isotonic, fitted on validation rows only; selection uses the
raw score, because isotonic regression collapses $\sim$800k distinct scores onto $\sim$100
levels and a quantile cut on the calibrated column silently drops part of the seed set.
% source: docs/paper/evaluation_protocol.md §5
Neighbourhood sampling is causal at the level of batches, not of single edges. The library
used here cannot express edge-level temporal sampling, so a fallback loader cuts batches by size
or by 1{,}800 seconds of seed time, whichever comes first, and samples neighbours up to the
batch's latest seed. A scored edge can therefore see neighbour edges, and other edges of its own
batch, up to 1{,}800\,s later than itself. The measured look-ahead has median 180\,s (95th
percentile 720\,s) at the scoring batch of 2048 used for MEGA-GNN and MEGA-GenAgg, and median
600\,s (95th percentile 1{,}440\,s) at the batch of 8192 used for the shipped Multi-GINe scores.
It applies to all three neural detectors.
% source: results/notes/p7_batch_span.json; results/notes/p5_required_changes.md (R-LA)

\subsection{Statistics}
\label{sec:data-stats}

Intervals are 2.5/97.5 percentile bootstrap intervals over \emph{truth rings} with detector seeds
averaged inside each resample; a ring drawn $k$ times counts $k$ times. The pre-registered matched-ring tests are paired
sign-flip permutations, Holm-corrected~\citep{holm1979}. They are \textbf{blind to precision}: a
method can emit far more rings, match about as many and test as indistinguishable (the full
pipeline on GFP-XGBoost scores sits 13.9 ring-F1 points below Multi-GINe's at
$p_{\mathrm{Holm}}=0.807$). We therefore add a \textbf{precision-aware paired bootstrap} that
resamples truth rings and off-ring clusters of predictions together and evaluates both arms on
the same draw, so ring precision responds to false rings. We also report the learned stages
with inputs \emph{cross-fitted by ring}, so that fold models never score a ring they were trained
on, beside the day-folded inputs first used (Section~\ref{sec:limitations}). Details are in docs/PROTOCOL.md.
% source: results/notes/interval_consistency.md; evaluation/ring_stats.py (bootstrap_over_rings, paired_unit_bootstrap); docs/FINDINGS.md (section B preamble); results/notes/p2_ringfold.md

\paragraph{The protocol as a scorer.}
We release the protocol as a standalone scorer that takes any file of predicted rings and
reports every scope convention, the threshold sweep, observability counts and paired
comparisons; on our runs it reproduces the point estimates of every scope table in this paper
exactly (\url{https://anonymous.4open.science/r/<ID>}).
% source: results/notes/p1_scorer.md (90 cells, maximum absolute deviation 0.0); docs/PROTOCOL.md
